\documentclass[11pt]{article}
\usepackage{amsmath,amsfonts,amssymb,amsthm}
\usepackage{graphicx}
\usepackage{algorithm}
\usepackage{algorithmic}
\usepackage{float}
\usepackage{hyperref}

\title{High-precision randomized iterative methods \\ for the random feature method}
\author{Jingrun Chen and Longze Tan}
\date{}

\begin{document}
\maketitle

\section{Introduction}

The Random Feature Method (RFM) has emerged as an effective approach for solving partial differential equations (PDEs), bridging traditional numerical methods and machine learning-based algorithms. While RFM offers significant advantages in handling complex geometries with exponential convergence rates, its practical implementation faces a critical challenge: solving large-scale, ill-conditioned, and overdetermined sparse least squares problems.

Traditional solvers for these problems fall into two categories: direct methods (such as SVD and QR decomposition) that achieve high accuracy but with computational complexity of $O(mn\min(m,n))$, and iterative methods with per-iteration complexity of $O(mn)$. For RFM applications, where $m \gg n$ by several orders of magnitude, direct methods become computationally expensive, while classical iterative methods often fail to achieve high precision due to the ill-conditioning.

This paper addresses this challenge by introducing two high-precision randomized iterative methods: Count Sketch QR Preconditioned LSQR (CSQRP-LSQR) and Count Sketch SVD Preconditioned LSQR (CSSVDP-LSQR). These methods combine randomized sketching techniques with preconditioning strategies to transform ill-conditioned problems into well-conditioned ones, enabling efficient and accurate solutions.

\section{Algorithm Details}

\subsection{Core Components}

Both proposed methods are built upon three key components:

\begin{enumerate}
    \item \textbf{Count Sketch Technique}: A dimensionality reduction method that preserves the essential information of the original matrix while significantly reducing computational requirements.
    \item \textbf{Preconditioner Construction}: Using the sketched matrix to construct a preconditioner through either QR factorization or singular value decomposition (SVD).
    \item \textbf{Iterative Solver}: Applying the LSQR method to the preconditioned system.
\end{enumerate}

\subsection{Count Sketch Transform}

The count sketch transform is defined as $S = \Phi D \in \mathbb{R}^{s \times m}$, where:
\begin{itemize}
    \item $D$ is an $m \times m$ random diagonal matrix with each diagonal entry independently chosen to be +1 or -1 with equal probability
    \item $\Phi \in \{0,1\}^{s \times m}$ is a binary matrix with $\Phi_{h(i),i} = 1$ and all remaining entries 0
    \item $h: [m] \rightarrow [s]$ is a random map such that for each $i \in [m]$, $h(i) = j$ with probability $1/s$ for each $j \in [s]$
\end{itemize}

This technique provides an $\ell_2$-subspace embedding for the column space of a matrix $A$, meaning for all $x \in \mathbb{R}^n$:
\begin{equation}
(1-\epsilon)\|Ax\|^2_2 \leq \|SAx\|^2_2 \leq (1+\epsilon)\|Ax\|^2_2
\end{equation}

The count sketch method offers optimal computational complexity of $O(\text{nnz}(A))$ (where nnz denotes the number of non-zero elements) and requires only one non-zero entry per column, preserving the sparsity structure of the original matrix.

\subsection{CSQRP-LSQR Algorithm}

The CSQRP-LSQR method proceeds as follows:

\begin{enumerate}
    \item Choose an oversampling factor $\gamma > 1$ and set $s = \gamma n$ (with $s < m$)
    \item Generate an $s \times m$ count sketch matrix $S$ and compute $\tilde{A} = SA$
    \item Apply QR factorization to $\tilde{A}$ to obtain the upper triangular matrix $R$
    \item Compute the preconditioned matrix $B = AR^{-1}$
    \item Use the LSQR method to find an approximate solution $y^{(k)}$ of $\min_{y \in \mathbb{R}^n} \|By - b\|_2$
    \item Return $x^{(k)} = R^{-1}y^{(k)}$
\end{enumerate}

\subsection{CSSVDP-LSQR Algorithm}

The CSSVDP-LSQR method is similar but uses SVD instead of QR factorization:

\begin{enumerate}
    \item Choose an oversampling factor $\gamma > 1$ and set $s = \gamma n$ (with $s < m$)
    \item Generate an $s \times m$ count sketch matrix $S$ and compute $\tilde{A} = SA$
    \item Compute the compact SVD $U\Sigma V^T$ of matrix $\tilde{A}$, where $r = \text{rank}(\tilde{A})$, $U \in \mathbb{R}^{s \times r}$, $\Sigma \in \mathbb{R}^{r \times r}$, $V \in \mathbb{R}^{n \times r}$, and let $P = V\Sigma^{-1}$
    \item Compute the preconditioned matrix $B = AP$
    \item Use the LSQR method to compute an approximate solution $y^{(k)}$ for $\min_{y \in \mathbb{R}^n} \|By - b\|_2$
    \item Return $x^{(k)} = Py^{(k)}$
\end{enumerate}

The CSSVDP-LSQR method is particularly effective for matrices that are rank-deficient or have very high condition numbers, as it performs a truncated SVD that effectively filters out numerically insignificant singular values.

\section{Theoretical Findings}

\subsection{Condition Number Bounds}

A key theoretical finding is that the condition number of the preconditioned matrix is bounded independently of the condition number of the original matrix:

\begin{theorem}
Let matrix $A$ in the least squares problem be of full column rank and the construction of the count sketch matrix $S$ satisfy certain conditions. Then, with probability at least $1-\delta$:
\begin{itemize}
    \item The upper triangular matrix $R$ obtained from the QR decomposition of $\tilde{A}$ is invertible
    \item The condition number of the preconditioned matrix $AR^{-1}$ satisfies $\kappa(AR^{-1}) \leq \sqrt{\frac{1+\epsilon}{1-\epsilon}}$
    \item The preconditioned least squares problem has a unique solution
\end{itemize}
\end{theorem}

This result shows that the proposed method can transform an ill-conditioned problem into a well-conditioned one with high probability, with the new condition number bounded by a constant that depends only on $\epsilon$.

\subsection{Error Estimates}

For the CSQRP-LSQR method, the paper provides detailed error estimates:

\begin{theorem}
Under the same conditions as above, and letting $x^*$ and $y^*$ be the least squares solutions to the original and preconditioned problems respectively, given a tolerance $\tau > 0$, when the LSQR method is used to solve the preconditioned system, if the number of iteration steps satisfies:
\begin{equation}
k \geq \frac{\ln(2) + |\ln(\tau)|}{|\ln(\epsilon)|}
\end{equation}
then:
\begin{equation}
\|y^{(k)} - y^*\|_{B^TB} \leq \tau \|y^*\|_{B^TB}
\end{equation}

This implies that the error in the original problem satisfies:
\begin{equation}
\|x^{(k)} - x^*\|_{A^TA} \leq \tau \|x^*\|_{A^TA}
\end{equation}
\end{theorem}

The paper also provides error estimates for the CSSVDP-LSQR method when applied to approximately rank-deficient problems, showing that under certain conditions, the method can determine the effective rank correctly and achieve high precision.

\subsection{Convergence Rate Analysis}

The paper shows that the convergence rate of LSQR on the preconditioned system is determined by:
\begin{equation}
\frac{\|y^{(k)} - y^*\|_{B^TB}}{\|y^{(0)} - y^*\|_{B^TB}} \leq 2\left(\frac{\sqrt{\kappa(B^TB)} - 1}{\sqrt{\kappa(B^TB)} + 1}\right)^k
\end{equation}

Since $\kappa(B^TB) \leq \frac{1+\epsilon}{1-\epsilon}$, this leads to exponential convergence with rate bounded by $\epsilon$, enabling high precision with relatively few iterations.

\section{Numerical Examples and Results}

\subsection{Two-Dimensional PDEs}

The paper presents extensive numerical results for various PDEs, including:

\begin{itemize}
    \item \textbf{Timoshenko Beam Problem}: Comparing various methods, CSSVDP-LSQR achieved relative errors of $10^{-17}$ while CSQRP-LSQR reached $10^{-13}$. In contrast, standard iterative methods like LSQR and LSMR could only achieve $10^{-4}$ after 10,000 iterations.
    
    \item \textbf{Stokes Flow Problem}: For a square domain with three holes, CSSVDP-LSQR achieved errors of $10^{-8}$ in $u$, $10^{-9}$ in $v$, and $10^{-5}$ in pressure $p$, significantly better than other methods.
    
    \item \textbf{Elasticity Problem with Complex Geometry}: Both CSQRP-LSQR and CSSVDP-LSQR methods achieved relative errors around $10^{-11}$ with errors in physical variables around $10^{-5}$.
\end{itemize}

A key finding was that using a count sketch matrix with $s = 3n$ (i.e., $\gamma = 3$) generally yielded optimal total time and accuracy across problems, despite theoretical bounds suggesting $s = O(n^2)$.

\subsection{Three-Dimensional PDEs}

For three-dimensional problems, the methods demonstrated similar effectiveness:

\begin{itemize}
    \item \textbf{Poisson Equation}: Both CSQRP-LSQR and CSSVDP-LSQR achieved relative errors down to $10^{-14}$ as problem size increased.
    
    \item \textbf{Helmholtz Equation}: For low-frequency solutions, CSQRP-LSQR achieved errors down to $10^{-14}$, while for high and mixed-frequency problems, it reached $10^{-12}$.
    
    \item \textbf{Elasticity Problem}: CSQRP-LSQR achieved relative $L^2$ errors down to $10^{-15}$, outperforming direct methods in computational time for large problems.
    
    \item \textbf{Stokes Flow}: For a cubic domain with inhomogeneous boundary conditions, CSQRP-LSQR reached errors of $10^{-16}$ in 63 iterations.
\end{itemize}

\subsection{Problems with Rank Deficiency or Infinite Condition Numbers}

The paper evaluated CSSVDP-LSQR on matrices from the Florida Sparse Matrix Collection with rank deficiency or infinite condition numbers:

\begin{itemize}
    \item CSSVDP-LSQR consistently achieved relative errors around $10^{-15}$ in 24-40 iterations.
    
    \item The condition number of preconditioned matrices $AP$ ranged from 3 to 6, compared to original condition numbers reaching $10^{43}$ or infinity.
    
    \item CSQRP-LSQR was less effective for rank-deficient problems but still worked for some cases with condition numbers as high as $10^{15}$.
\end{itemize}

The paper observed that singular values of preconditioned matrices were more closely clustered near 1, explaining the improved convergence behavior.

\section{Implications and Advantages}

\subsection{Computational Efficiency}

The proposed methods offer several computational advantages:

\begin{itemize}
    \item CSQRP-LSQR is computationally faster on large-scale problems compared to the state-of-the-art sparse direct solver SPQR.
    
    \item The count sketch technique preserves sparsity patterns, allowing efficient operations without explicitly forming the sketching matrix.
    
    \item As problem size increases, the advantage of these methods over direct solvers becomes more pronounced.
\end{itemize}

\subsection{Accuracy and Reliability}

The accuracy achieved is comparable to direct methods:

\begin{itemize}
    \item The relative least squares error precision matches that of SPQR and Householder QR methods.
    
    \item The methods handle problems with condition numbers up to $10^{20}$ while maintaining high precision.
    
    \item CSSVDP-LSQR successfully solves problems with rank deficiency or infinite condition numbers where direct methods fail.
\end{itemize}

\subsection{Practical Implications for RFM}

These methods significantly enhance the applicability of RFM for solving PDEs:

\begin{itemize}
    \item They enable high-precision solutions necessary for RFM to achieve its theoretical exponential convergence rate.
    
    \item They make RFM practical for complex geometries and high-complexity solutions.
    
    \item The enhanced solver performance allows RFM to compete with both traditional numerical methods and machine learning approaches.
\end{itemize}

The paper notes that while theory suggests a sketching dimension of $s = O(n^2)$ is required, numerical experiments show that $s = 3n$ works well in practice, significantly improving the practical utility of count sketch transformations.

\subsection{Future Research Directions}

The work opens several promising research directions:

\begin{itemize}
    \item Closing the gap between theoretical requirements ($s = O(n^2)$) and practical performance ($s = O(n)$) for count sketch dimensions.
    
    \item Further optimizing the preconditioning strategies for specific problem structures.
    
    \item Exploring parallel implementations to further enhance computational efficiency.
    
    \item Applying the methods to even broader classes of PDEs and other scientific computing problems.
\end{itemize}

\end{document} 